\documentclass[preprint,12pt]{elsarticle}

\usepackage{amsmath,amssymb,amsthm}
\usepackage{booktabs}
\usepackage{tabularx}
\usepackage{graphicx}
\usepackage{microtype}
\usepackage{xcolor}
\usepackage{placeins}
\usepackage{hyperref}
\hypersetup{hidelinks}

\journal{Journal of Information Security and Applications}

\newcommand{\Lang}{\mathcal{L}}
\newcommand{\Rank}{\mathsf{Rank}}
\newcommand{\Unrank}{\mathsf{Unrank}}
\newcommand{\Perm}{\Pi}
\newcommand{\CommitEvidence}{\mathsf{CommitEvidence}}
\newtheorem{theorem}{Theorem}
\newtheorem{lemma}{Lemma}
\newtheorem{definition}{Definition}
\newcolumntype{L}[1]{>{\raggedright\arraybackslash}p{#1}}
\newcommand{\CorpusTotal}{270}
\newcommand{\CorpusTranslated}{121}
\newcommand{\CorpusUnsupported}{149}
\newcommand{\CorpusCompiled}{121}
\newcommand{\CorpusTimeout}{0}
\newcommand{\CorpusMedianStates}{65}
\newcommand{\CorpusNinetyFifthStates}{508}
\newcommand{\CorpusMaxStates}{49401}
\newcommand{\CorpusMedianPayloadMiB}{0.48}
\newcommand{\CorpusNinetyFifthPayloadMiB}{18.75}
\newcommand{\CorpusMaxPayloadMiB}{439.40}
\newcommand{\CorpusMedianRssMiB}{0.11}
\newcommand{\CorpusNinetyFifthRssMiB}{55.59}
\newcommand{\CorpusMaxRssMiB}{2064.57}
\newcommand{\CorpusMedianRankUs}{23.57}
\newcommand{\CorpusMedianUnrankUs}{33.19}
\newcommand{\WarmMedianUs}{128.90}
\newcommand{\WarmNinetyFifthUs}{242.54}
\newcommand{\ColdMedianMs}{5.53}
\newcommand{\SequenceCount}{100,000}
\newcommand{\SequenceDomain}{2,351,927,613,440,572,602,240}
\newcommand{\SequenceBits}{70.99}
\newcommand{\SequenceSeconds}{23.55}
\newcommand{\SequenceCalls}{200,000}
\newcommand{\PermutationEligible}{97}
\newcommand{\PermutationIneligible}{24}
\newcommand{\PermutationSamples}{3,104}
\newcommand{\PermutationMedianWalks}{1}
\newcommand{\PermutationNinetyFifthWalks}{3}
\newcommand{\PermutationMaxWalks}{9}
\newcommand{\PermutationCapHits}{0}
\newcommand{\RotationCases}{30}
\newcommand{\RotationCommitted}{15}
\newcommand{\RotationUnknown}{11}
\newcommand{\RotationPrepared}{2}
\newcommand{\RotationAborted}{2}
\newcommand{\DistributionSamples}{100,000}
\newcommand{\DichopileTvd}{0.00399}
\newcommand{\EpscdTvd}{0.00234}
\newcommand{\SparseMedianRetries}{6,696}
\newcommand{\SparseNinetyNinthNineRetries}{75,963}
\newcommand{\TlaDistinctStates}{1,069}


\begin{document}

\begin{frontmatter}

\title{Exact Policy-Space Credential Derivation for Failure-Safe Legacy Password Rotation}

\author{Yuanyi Zhang}
\address{Independent Researcher}

\begin{abstract}
Enterprise networks still contain applications, appliances, and directories
whose only modifiable authentication secret is a policy-constrained password.
Replacing these systems is preferable, but during migration a manager must
generate reconstructible credentials and rotate them through remote interfaces
whose success is not always observable. We present Exact Policy-Space
Credential Derivation (EPSCD), a construction that compiles a bounded password
policy into a finite accepted language, computes its exact cardinality, and
builds inverse ranking functions. A context-separated key controls a
permutation of the exact rank domain; generation $g$ is unranked to the
credential at that position. The construction specializes established
finite-language and rank-then-encipher techniques, giving exact compliance,
deterministic reconstruction, exact $\log_2|\Lang_P|$, and strict
non-repetition within one credential lineage. We then couple the sequence to an
evidence-bounded rotation protocol. Candidate metadata is durably prepared
before submission, and the committed generation advances only when an
adapter-specific predicate establishes acceptance of that candidate. Timeouts,
lost responses, unsafe verification, and contradictory observations preserve
both old and new reconstructability in an unknown-outcome state. In a clean run
over the published \CorpusTotal-policy SOUPS corpus, \CorpusCompiled\ of
\CorpusTranslated\ exactly translated policies compile under a fixed budget.
A \SequenceCount-generation run has no policy
violation, duplicate, or replay mismatch. \RotationCases\ fault scenarios across HTTP
form and LDAP-style adapter models preserve the two rotation invariants. A
bounded TLA+ model explores \TlaDistinctStates\ distinct states without invariant violations,
while two negative controls are detected. The artifact is a research prototype:
its FF1 backend uses bounded cycle walking, its adapter experiments are local
semantic models, and it does not make weak legacy policies strong.
\end{abstract}

\begin{keyword}
deterministic password generation \sep password policy \sep regular language
\sep format-preserving encryption \sep credential rotation \sep fault recovery
\end{keyword}

\end{frontmatter}

\section{Introduction}

Modern authentication should use phishing-resistant mechanisms when the
relying system can be changed. In practice, enterprise networks continue to
contain vendor applications, appliances, directories, and internally developed
services whose available interface is a password field. Passwords therefore
remain a migration and compatibility problem even when an organization has a
passwordless target architecture~\cite{herley2012persistence}. A management
client for these targets must answer two different questions. Which password
does generation $g$ denote under the target policy? When a remote change is
attempted, what evidence justifies changing the locally committed generation?

Published deterministic generators demonstrate site-specific derivation
without changing the target verifier. PwdHash and Password Multiplier derive a
site credential from a master secret and site context
\cite{ross2005pwdhash,halderman2005multiplier}; PALPAS separates a high-entropy
secret from non-secret per-service state and policy information~\cite{palpas}.
Published policy languages and verified generators separately establish that
machine-readable constraints and uniform compliant generation are feasible
\cite{gautam2022pcp,grilo2022verified}. These are essential precedents, but they
do not define a generation-indexed, without-replacement credential sequence
coupled to a failure-aware remote password-change journal.

Naively hashing $(K,\mathit{service},g)$ and reducing modulo a legal-space size
has two avoidable defects. Reduction can bias the marginal distribution, and
independent evaluations can collide at two generations. Constructing required
character classes and shuffling with unbiased indices also does not generally
make the final distribution uniform over the complete accepted set. The
correct abstraction is the finite language accepted by the entire policy.
Counting and ranking finite languages, arbitrary-domain ciphers, and
rank-then-encipher format-preserving encryption are established results
\cite{goldberg1985ranking,black2002finite,bellare2009fpe,oudinet2013uniform}.
Our method does not rename these results as new mathematics; it binds them to a
credential lineage and generation contract.

Remote rotation presents a separate systems failure. An HTTP success page, a
completed TCP exchange, or a submitted LDAP request need not prove that the
new credential is the durable remote state. Conversely, a server may commit an
update whose response is lost. Advancing local generation on a transport event
can make the manager reconstruct a credential that the target rejects. Rolling
back after every timeout can lose a credential that the target accepted. The
safe result is sometimes explicitly unknown.

This paper makes exactly two technical contributions:

\begin{enumerate}
\item \textbf{Exact Policy-Space Credential Sequence.} We specify a bounded
  policy compiler with exact arbitrary-precision suffix counts, inverse
  $\Rank/\Unrank$ functions, an explicit credential lineage, and an audited
  keyed permutation over $[0,N_P)$. The sequence provides exact compliance,
  deterministic reconstruction, exact effective bits, and strict same-lineage
  non-repetition. Cross-lineage equality is quantified rather than hidden by a
  password-history mechanism.
\item \textbf{Evidence-Bounded Failure-Safe Legacy Rotation.} We define a
  durable prepare--submit--verify protocol and an adapter capability contract.
  Local generation advances only after adapter-defined evidence supports the
  candidate credential; uncertain outcomes retain both deterministic
  reconstruction paths and permit later reconciliation.
\end{enumerate}

The contributions are a credential construction and a lifecycle protocol, not
a new finite-language algorithm, cryptographic primitive, distributed
transaction primitive, or authentication mechanism. EPSCD is intended for
legacy compatibility while a target is migrated to modern authentication.

\section{Related work and novelty boundary}

\subsection{Deterministic and policy-aware credential generation}

PwdHash derives site-separated passwords inside a browser extension without
server modification~\cite{ross2005pwdhash}. Password Multiplier uses a similar
master-secret/site-name model and its evaluation discusses changing the site
argument for rotation~\cite{halderman2005multiplier}. PALPAS derives passwords
from a high-entropy secret and per-service salts while synchronizing non-secret
parameters~\cite{palpas}. These works establish deterministic legacy-verifier
compatibility and service separation as prior art.

Horsch et al. define a password-requirements markup language
\cite{horsch2016prml}. Gautam et al. design a composition-policy description
language and release the 270-policy workbook used in our evaluation
\cite{gautam2022pcp}. Grilo et al. verify policy compliance and uniform random
generation in EasyCrypt and integrate the generator into Bitwarden
\cite{grilo2022verified}. EPSCD neither claims that policy representation nor
uniform compliant random generation is new. Its different object is a keyed,
deterministic, injective sequence indexed by credential generation.

\subsection{Finite-language counting and exact-domain permutations}

Uniform generation from formal languages and language ranking have a long
literature~\cite{hickey1983uniform,goldwurm1995random,denise1999uniform,
bernardi2012linear,oudinet2013uniform,goldberg1985ranking}. We use exact suffix
counts and lexicographic unranking, standard dynamic-programming techniques.
Dichopile~\cite{oudinet2013uniform} is the published algorithmic comparison in
our distribution and cost experiments. The algorithms answer different
questions: Dichopile draws independent random words, whereas EPSCD assigns a
stable permutation position to each credential generation.

Black and Rogaway study ciphers on arbitrary finite domains
\cite{black2002finite}. Bellare et al. formalize FPE, analyze cycle walking, and
describe rank-then-encipher for complex formats including regular languages
\cite{bellare2009fpe}. Consequently, the mapping
$\Unrank(\Perm(\Rank(x)))$ is not our invention. EPSCD specializes the approach
to credential metadata and rotation. The prototype reuses FF1 and follows the
one-million-element minimum-domain direction in the NIST revision draft
\cite{nist38g}; it does not propose a cipher.

\subsection{Key recovery and operational secret handling}

Threshold recovery and multi-factor key derivation are established topics
\cite{shamir1979share,nair2023mfkdf,scarlata2024mfkdf}. They are operational
dependencies rather than EPSCD contributions. In the intended profile, a
complete Root Key is not persistently stored. A managed endpoint retains one
Shamir share and removable media retains a second share under password-based
authenticated encryption; the password protects the share envelope and does
not generate the share. A required operation reconstructs the Root Key
transiently, derives the credential key and output, and clears temporary
material on a best-effort basis. A third network-resident share is reserved for
alternative recovery. This profile stores neither service-password values nor
a complete Root Key, but it cannot protect plaintext during a legitimate
operation on an already compromised endpoint.

\subsection{Published-property comparison}

Table~\ref{tab:prior-art} reports only properties specified by published work.
``Not specified'' is not a claim that an implementation cannot be extended.

\begin{table*}[t]
\centering
\footnotesize
\setlength{\tabcolsep}{3pt}
\caption{Published work and the two EPSCD research properties.}
\label{tab:prior-art}
\begin{tabularx}{\textwidth}{L{2.8cm}XXXX}
\toprule
Work & Deterministic service credential & Exact accepted-space cardinality & Same-lineage permuted sequence & Evidence-bounded remote rotation \\
\midrule
PwdHash~\cite{ross2005pwdhash} & Yes & Not specified & Not specified & Not specified \\
Password Multiplier~\cite{halderman2005multiplier} & Yes & Not specified & Not specified & Not specified \\
PALPAS~\cite{palpas} & Yes & Not specified & Not specified & Not specified \\
Verified password generation~\cite{grilo2022verified} & Random generation & Policy semantics verified & No & No \\
Dichopile~\cite{oudinet2013uniform} & No & Exact counting in our reproduction & No & No \\
EPSCD & Yes & Yes & Yes & Yes \\
\bottomrule
\end{tabularx}
\end{table*}

\section{System and threat model}

\subsection{Entities and persistent state}

The client manages a legacy target, a service/account identity, an authenticated
policy description, and Root Key material. The target exposes a password-change
adapter and possibly a separate verification or version-readback interface. It
does not participate in EPSCD and need not store EPSCD metadata.

The baseline persists identifiers, a random credential salt, an explicit
lineage identifier, policy identity/version/epoch/hash, root generation,
committed credential generation, adapter metadata, and a pending operation
journal. It does not persist the current or previous service-password values.
During display, submission, and verification, a derived password necessarily
exists transiently.

\subsection{Adversary and assumptions}

The adversary may read public policy and generation metadata; observe known
$(g,\mathit{pwd}_g)$ pairs; steal the metadata database; drop, delay, duplicate,
or reorder network traffic; replay a stale local snapshot; and crash the client
at a persistence or network boundary. An adapter can be unavailable or return
ambiguous observations. We separately discuss adapter compromise.

The authenticated journal is assumed to preserve integrity on the current
device. Detecting rollback of the journal itself requires an external or
non-rollbackable freshness checkpoint. The selected permutation is assumed to
meet its stated PRP security goal on the used domains. Policy and generation
metadata are not secret. A fully compromised endpoint during a legitimate
derivation can capture the Root Key, credential key, or plaintext password.

\subsection{Goals and non-goals}

The goals are exact policy membership, deterministic reconstruction,
same-lineage non-repetition, exact domain measurement, and safe local progress
under remote ambiguity. Non-goals include forward secrecy after credential-key
compromise, global non-repetition across independent lineages, availability
against a malicious operator, complete constant-time behavior, and making a
small accepted language resistant to exhaustive search.

\begin{figure*}[t]
\centering
\small
\setlength{\tabcolsep}{5pt}
\resizebox{\textwidth}{!}{%
\begin{tabular}{ccccccccc}
\fbox{Policy $P$} & $\rightarrow$ & \fbox{Canonical IR} & $\rightarrow$ & \fbox{Finite automaton} & $\rightarrow$ & \fbox{Exact counts} & $\rightarrow$ & $N_P,\Rank_P,\Unrank_P$ \\
&&&&&&&& $\downarrow$ \\
\fbox{$K_{root}$ + context $C$} & $\rightarrow$ & \fbox{$K_{cred}$ + tweak $T$} & $\rightarrow$ & \fbox{generation $g$} & $\rightarrow$ & \fbox{$r=\Perm_{K,T}(g)$} & $\rightarrow$ & \fbox{$\Unrank_P(r)$} \\
\end{tabular}
}
\caption{EPSCD data flow. Policy compilation defines one exact finite domain; credential generation is the input to one context-bound permutation, not part of its tweak.}
\label{fig:pipeline}
\end{figure*}


\section{Exact Policy-Space Credential Sequence}

\subsection{Bounded policy language}

Let $P$ be a canonical bounded policy over alphabet $\Sigma$, with lengths in
$[L_{\min},L_{\max}]$. The implemented fragment supports allowed and forbidden
characters; minimum and maximum class counts; fixed positions, prefix, and
suffix; forbidden first/last sets; a per-character maximum; identical and
sequential run limits; and finite forbidden substrings. Their product state is
a finite deterministic transition system

\[
 A_P=(Q,\Sigma,\delta,q_0,F).
\]

The exact accepted language is

\[
 \Lang_P=\{w\in\Sigma^*:L_{\min}\le |w|\le L_{\max},
                         \delta^*(q_0,w)\in F\}.
\]

Unsupported source semantics are rejected, not approximated. The compiler also
enforces a state budget because the product of class counters, substring state,
run state, and position constraints can be exponential in the policy
description.

\subsection{Exact suffix counts}

For position $i$, target length $\ell$, and automaton state $q$, define

\[
C_\ell(i,q)=
\begin{cases}
1,& i=\ell\ \land\ q\in F,\\
0,& i=\ell\ \land\ q\notin F,\\
\sum_{c\in\Sigma:\delta(q,c)\downarrow}
  C_\ell(i+1,\delta(q,c)),& i<\ell.
\end{cases}
\]

All counts use arbitrary-precision integers. The exact domain size is

\[
 N_P=|\Lang_P|=\sum_{\ell=L_{\min}}^{L_{\max}}C_\ell(0,q_0),
 \qquad H_{\mathrm{eff}}=\log_2 N_P.
\]

For $Q=|Q|$, maximum length $L$, and alphabet size $s$, a dense suffix table
uses $O(LQ)$ big integers and $O(LQs)$ transition/count operations after the
automaton is constructed. Big-integer bit complexity depends on $\log N_P$.
The prototype reports state count, count payload, process RSS, compile time,
and online rank/unrank latency rather than hiding these costs in one average.

\subsection{Ranking and unranking}

Fix a canonical order on lengths and on $\Sigma$. Unranking first selects the
length interval containing rank $r$. At each position, it considers outgoing
characters in canonical order; the completion count following each edge is an
interval width. It selects the interval containing $r$ and subtracts preceding
widths. Ranking applies the inverse accumulation to an accepted word.

\begin{theorem}[Bijection]
For a successfully compiled non-empty policy, $\Rank_P:\Lang_P\rightarrow
[0,N_P)$ and $\Unrank_P:[0,N_P)\rightarrow\Lang_P$ are mutual inverses.
\end{theorem}

\begin{proof}
At each prefix, outgoing edge intervals are disjoint, adjacent, and their widths
sum to the suffix count of the current state by the recurrence. Therefore every
rank selects exactly one edge and every accepted prefix contributes exactly the
sum of intervals preceding its edge. Induction over remaining positions gives
$\Rank_P(\Unrank_P(r))=r$. Applying the same argument to the unique path of an
accepted word gives $\Unrank_P(\Rank_P(w))=w$.
\end{proof}

The decoder performs no random candidate rejection. This statement does not
apply to the concrete permutation layer.

\subsection{Context-separated exact-domain sequence}

Scheme v2 persists a public random $\mathit{lineageID}$. A credential key is
derived using canonical length-delimited metadata:

\begin{align*}
K_{\mathrm{cred}}={}&\mathrm{HKDF}(K_{\mathrm{root}},
  \texttt{EPSCD/credential-key/scheme-v2}\parallel\\
 &\mathit{vaultID}\parallel\mathit{serviceID}\parallel\mathit{accountID}
 \parallel\mathit{lineageID}\parallel\mathit{credentialSalt}\parallel\\
 &\mathit{rootGeneration}\parallel\mathit{policyVersion}
 \parallel\mathit{policyEpoch}).
\end{align*}

The tweak additionally binds the policy identifier and canonical policy hash.
Generation is deliberately the permutation input, not part of the tweak:

\[
 r_g=\Perm_{K_{\mathrm{cred}},T}(g),\qquad
 \mathit{pwd}_g=\Unrank_P(r_g),\qquad 0\le g<N_P.
\]

The concrete backend applies AES-256 FF1 to the smallest binary domain
$[0,M)$ containing $[0,N_P)$ and cycle-walks until the result is below $N_P$.
It fails closed after 1,024 walks, supports domains of at least one million
elements and at most 512 bits, and never falls back to modulo reduction. Since
$M<2N_P$, the accepted density is greater than one half for this selected
superset. The implementation is not claimed to be constant time.

\begin{theorem}[Exact compliance and deterministic reconstruction]
Every successful derivation returns an element of $\Lang_P$. Identical root,
canonical context, policy, generation, and backend produce the same password.
\end{theorem}

\begin{proof}
The backend returns $r_g\in[0,N_P)$ and Theorem~1 establishes that
$\Unrank_P(r_g)\in\Lang_P$. All transformations are deterministic for fixed
inputs.
\end{proof}

\begin{theorem}[Same-lineage non-repetition]
For fixed key, tweak, and policy, successful derivations at $g_1\ne g_2$ return
different passwords until the domain is exhausted.
\end{theorem}

\begin{proof}
$\Perm_{K,T}$ is injective and $\Unrank_P$ is injective. Their composition is
injective.
\end{proof}

\begin{theorem}[Ideal-permutation marginal]
For a fixed generation and a uniformly selected ideal random permutation key,
each $w\in\Lang_P$ occurs with probability $1/N_P$.
\end{theorem}

\begin{proof}
The image of a fixed input under an ideal random permutation is uniform on
$[0,N_P)$. Bijection under $\Unrank_P$ preserves the mass of every element.
\end{proof}

For the concrete backend, Theorem~4 is replaced by the applicable computational
PRP claim. Known $(g,\mathit{pwd}_g)$ pairs are known input/output pairs after
ranking and do not provide forward secrecy.

\subsection{Lineage change and collision bound}

Rekeying, root replacement, or a semantically changed policy starts a new
lineage. Strict non-repetition is intentionally scoped to one lineage. Let
$H$ be historical passwords, $\Lang_{\mathrm{new}}$ the new language,
$h=|H\cap\Lang_{\mathrm{new}}|$, and $N=|\Lang_{\mathrm{new}}|$. Under an
independent ideal permutation, the next output collides with this relevant
history with probability $h/N$. For $m$ future outputs, a simple union bound is

\[
 \Pr[\mathrm{collision}]\le \min\left(1,\frac{hm}{N}\right).
\]

The prototype returns this bound as an exact rational. It stores no password
history in the baseline. A deployment can configure a minimum
$H_{\mathrm{eff}}$ and reject or warn on a weak legacy policy.

\section{Evidence-Bounded Failure-Safe Rotation}

\subsection{Why remote success is an evidence question}

A remote password-change operation crosses independent persistence and network
boundaries. The request may never arrive; the server may commit and lose its
response; a proxy may return a success page before durable update; replicated
authentication nodes may temporarily disagree; or verification may be unsafe
because failed attempts consume a lockout budget. Therefore no universal
mapping from transport status to committed credential exists.

\begin{figure*}[t]
\centering
\resizebox{\textwidth}{!}{$
\fbox{\textsc{Active}$(g)$}
\xrightarrow{\text{persist }(g,g+1,opID)}
\fbox{\textsc{Prepared}}
\xrightarrow{\text{submit}}
\fbox{\textsc{Submitted}}
\xrightarrow{\text{adapter observations}}
\fbox{\textsc{Verifying}}
$}

\vspace{6pt}
\scriptsize
\begin{tabularx}{0.96\textwidth}{>{\centering\arraybackslash}X>{\centering\arraybackslash}X>{\centering\arraybackslash}X}
$\mathsf{CommitEvidence}=\mathsf{true}$
$\rightarrow\textsc{Committed}(g+1)$ &
Conclusive old-authoritative evidence
$\rightarrow\textsc{Aborted}(g)$ &
Insufficient or contradictory evidence
$\rightarrow\textsc{UnknownOutcome}$ \\
\multicolumn{3}{c}{HTTP 200 $\not\Rightarrow$ commit; \quad
\textsc{UnknownOutcome} preserves reconstructibility of both $g$ and $g+1$.}
\end{tabularx}
\caption{Evidence-bounded remote rotation. Persisting a candidate and sending
an update never commits local metadata. The adapter contract determines whether
new acceptance, new-only acceptance, atomic evidence, or authoritative version
readback is sufficient.}
\label{fig:rotation-lifecycle}
\end{figure*}


\subsection{Adapter capability and evidence contract}

Each adapter declares whether it can verify the new credential, safely verify
the old credential, obtain atomic success evidence, read an authoritative
remote version, and attach an idempotency key. Its evidence requirement is one
of:

\begin{itemize}
\item \textsc{NewAcceptance}: a conclusive candidate probe is sufficient;
\item \textsc{NewOnly}: the candidate succeeds and a safe old probe
  conclusively fails;
\item \textsc{AuthoritativeVersion}: an expected remote version/readback is
  observed;
\item \textsc{UnknownOnly}: automatic commit evidence cannot be established.
\end{itemize}

An adapter-specific predicate
$\CommitEvidence(A,O)$ evaluates declared capabilities $A$ and persisted
observations $O$. A generic update success flag is recorded but ignored by the
predicate. The old probe is omitted when unsafe; this can narrow the conclusion
to new acceptance or force an unknown outcome rather than spending the lockout
budget.

LDAP provides a standardized password-modify operation~\cite{rfc3062}, but an
adapter must still define what a successful operation and subsequent readback
mean in its deployment. The protocol does not treat the existence of a standard
operation as universal evidence of replicated authentication state.

\subsection{Durable transition system}

Let $g$ be the committed generation. Before any possible network send, the
client persists

\[
\textsc{Prepared}(g,g+1,opID,\mathit{lineage},\mathit{policy},A).
\]

The record contains metadata sufficient to reconstruct both candidates, not
their password values. Submission enters \textsc{Submitted}; observations enter
\textsc{Verifying}. The only terminal decisions are:

\begin{align*}
\textsc{Committed}(g+1)&\quad\text{if }\CommitEvidence(A,O),\\
\textsc{Aborted}(g)&\quad\text{if evidence conclusively retains the old state},\\
\textsc{UnknownOutcome}&\quad\text{otherwise}.
\end{align*}

\textsc{UnknownOutcome} neither increments nor automatically rolls back. It
retains the base generation, candidate generation, operation identity, and
context needed to reconstruct both. Later reconciliation adds observations to
the journal. A stale operation identifier cannot commit a different candidate.

\begin{theorem}[Commit safety]
Assuming authenticated journal writes and correct adapter classification, the
committed generation cannot advance from $g$ to $g+1$ without persisted evidence
that satisfies the declared adapter predicate for the corresponding candidate.
\end{theorem}

\begin{proof}
The prepare, submit, verify, crash, unknown, and abort transitions do not modify
the committed generation. The only transition that does is commit. Its guard
requires the persisted candidate identity and a true value of the adapter
evidence predicate. Induction over a transition trace establishes the invariant.
\end{proof}

\begin{theorem}[Uncertainty preservation]
Every reachable \textnormal{\textsc{UnknownOutcome}} record retains sufficient authenticated
metadata to reconstruct both $\mathit{pwd}_g$ and $\mathit{pwd}_{g+1}$.
\end{theorem}

\begin{proof}
Preparation persists both generation contexts before submission. Every
transition into \textsc{UnknownOutcome} preserves that record; deletion of
either path is permitted only by a terminal commit or abort transition.
\end{proof}

These are safety results under the stated model, not claims of liveness or a new
distributed transaction primitive.

\section{Implementation}

The Rust prototype uses a canonical JSON policy IR, a reachable-state compiler,
and a dense arbitrary-precision suffix-count table. The implementation caps
reachable states at 250,000 and policy-domain length at the documented bound.
Rank and unrank scan canonical outgoing characters; policy acceptance reuses the
compiled transitions. HKDF-SHA-256 derives scheme-v2 credential keys. The
permutation dependency implements FF1 with AES-256; EPSCD supplies range,
domain, and cycle-walk guards.

Rotation metadata is held in authenticated credential-description records and a
SQLite store. The adapter contract is a typed structure and the evidence
classifier returns sufficient-new-acceptance, sufficient-new-only,
sufficient-atomic, sufficient-version, insufficient, or contradictory. The
local semantic adapters reopen an on-disk SQLite journal during crash cases and
contain no password column.

The artifact command runs Rust unit/property tests, the public-corpus compiler,
controlled-density and sequence experiments, rotation faults, and TLA+, then
generates JSON, CSV, and LaTeX tables. Paper numbers in this section are inputs
from those generated tables.

\section{Evaluation}

We ask six questions. RQ1 measures realistic policy coverage and compiler cost;
RQ2 checks exactness; RQ3 studies the cost of whole-candidate rejection as
acceptance density falls; RQ4 exercises a long deterministic sequence; RQ5
tests failure safety; and RQ6 checks whether different adapter capabilities fit
one lifecycle.

\subsection{Setup and measurement boundaries}

The policy corpus is the public historical workbook published with Gautam et
al.~\cite{gautam2022pcp}. It is a reproducibility corpus, not a claim about
current enterprise prevalence. Translation rejects unsupported semantics
instead of weakening them. Each translated policy receives fixed bounds of
250,000 states, 4~GiB RSS, and 60 seconds. Timing values are single-host
prototype measurements. The published Dichopile algorithm is reproduced with
exact BigUint arithmetic; we do not claim performance against the authors'
optimized experimental configuration.

\subsection{RQ1: policy coverage and exact-space cost}

\begin{table}[t]
\centering
\caption{Public policy corpus clean run.}
\label{tab:corpus}
\begin{tabular}{lr}
\toprule
Metric & Result \\
\midrule
Source policies & 270 \\
Exact translations & 121 \\
Compiled & 121 \\
Time limit & 0 \\
Median states & 65 \\
P95 states & 508 \\
Median compile (ms) & 42.73 \\
P95 compile (ms) & 694.99 \\
Median rank ($\mu$s) & 23.57 \\
Median unrank ($\mu$s) & 33.19 \\
\bottomrule
\end{tabular}

\end{table}

Table~\ref{tab:corpus} shows all \CorpusTotal\ source records. Of these,
\CorpusTranslated\ can be translated exactly into the implemented bounded IR
and \CorpusUnsupported\ are rejected for unsupported source semantics. Of the
attempted exact translations, \CorpusCompiled\ compile, with \CorpusTimeout\
reaching the fixed time limit. Among successful policies, median/P95 reachable
states are \CorpusMedianStates/\CorpusNinetyFifthStates\ and maximum is
\CorpusMaxStates. Median/P95 count payload is
\CorpusMedianPayloadMiB/\CorpusNinetyFifthPayloadMiB~MiB; maximum is
\CorpusMaxPayloadMiB~MiB. Sampled process RSS remains in the machine-readable
record, with the stated warning that short-lived peaks can be missed. These
tails justify reporting resource
limits rather than only a success percentage.

Median policy-level rank and unrank latencies are \CorpusMedianRankUs\ and
\CorpusMedianUnrankUs~$\mu$s.
On the separate 18-character performance policy, cached full derivation is
\WarmMedianUs~$\mu$s median and \WarmNinetyFifthUs~$\mu$s P95; cold compile-and-derive
is \ColdMedianMs~ms
median. Compiler cost is therefore paid at policy enrollment or change, not on
each routine credential reconstruction.

\subsection{RQ2: exactness and distribution checks}

Small-language tests enumerate every rank and verify acceptance. They check
\[
\Rank(\Unrank(r))=r
\quad\text{and}\quad
\Unrank(\Rank(w))=w.
\]
Property tests vary length and
class requirements. A count exceeding $2^{128}$ verifies that cardinality is
not truncated to a machine integer. Permutation tests check inversion and
injection on sampled inputs.

For the enumerable language $\{a,b0,\ldots,b9\}$, \DistributionSamples\ independent samples
from the exact Dichopile reproduction have empirical total-variation distance
\DichopileTvd\ from uniform; the test-only ideal-permutation EPSCD composition has
distance \EpscdTvd. These are finite-sample sanity checks, not proofs and not a
concrete FF1 security experiment.

\subsection{RQ3: controlled rejection density}

For independent whole-candidate rejection with acceptance density $\alpha$,
the retry count $R$ is geometric:

\[
 \mathbb{E}[R]=1/\alpha,\qquad
 \Pr[R>t]=(1-\alpha)^t.
\]

\begin{table}[t]
\centering
\caption{Whole-candidate rejection retries, 2,000 accepted outputs per density.}
\label{tab:density}
\begin{tabular}{rrrrrr}
\toprule
$\alpha$ & $1/\alpha$ & Median & P95 & P99 & P99.9 \\
\midrule
1e-01 & 10 & 7 & 28 & 40 & 60 \\
1e-02 & 100 & 69 & 291 & 449 & 933 \\
1e-03 & 1000 & 661 & 3012 & 4713 & 7427 \\
1e-04 & 10000 & 6696 & 31022 & 45236 & 75963 \\
\bottomrule
\end{tabular}

\end{table}

Table~\ref{tab:density} demonstrates the increasingly long online tail. At
$\alpha=10^{-4}$, the observed median is \SparseMedianRetries\ attempts and
P99.9 is \SparseNinetyNinthNineRetries.
EPSCD's language decoder follows one count-selected path and therefore has no
candidate-retry tail. This comparison isolates the mechanism; the selected FF1
backend has its own bounded cycle walks and must not be described as wholly
rejection-free.

Across the \CorpusCompiled\ compiled corpus policies, \PermutationEligible\
satisfy the FF1 backend's domain bounds and \PermutationIneligible\ are rejected
by that backend while remaining compiler successes. For \PermutationSamples\
measured generations, cycle walks have median \PermutationMedianWalks, P95
\PermutationNinetyFifthWalks, maximum \PermutationMaxWalks, and
\PermutationCapHits\ cap failures. Because the smallest binary superset has density above
one half, these observations are not the same experiment as the synthetic
$10^{-4}$ whole-candidate density.

\subsection{RQ4: sequence behavior}

\begin{table}[t]
\centering
\caption{Scheme-v2 long-sequence implementation check.}
\label{tab:sequence}
\begin{tabular}{lr}
\toprule
Metric & Result \\
\midrule
Generations & 100,000 \\
Policy violations & 0 \\
Duplicate passwords & 0 \\
Replay mismatches & 0 \\
Effective bits & 70.99 \\
Derivations/s & 8493 \\
\bottomrule
\end{tabular}

\end{table}

The run derives and replays \SequenceCount\ consecutive generations in a domain
of \SequenceDomain\ elements (\SequenceBits\ effective bits).
Table~\ref{tab:sequence} reports no policy violation, duplicate, or replay
mismatch. The run takes \SequenceSeconds\ seconds for \SequenceCalls\
derivations on the test host.
This validates the implementation path at scale; non-repetition itself follows
from Theorem~3 rather than empirical collision absence.

\subsection{RQ5--RQ6: faults and adapter diversity}

\begin{table}[t]
\centering
\caption{Fault matrix across two adapter semantics.}
\label{tab:rotation}
\begin{tabular}{lr}
\toprule
Metric & Result \\
\midrule
Adapters & 2 \\
Injected scenarios & 30 \\
Committed & 15 \\
Unknown outcome & 11 \\
Aborted & 2 \\
All safety checks & pass \\
\bottomrule
\end{tabular}

\end{table}

The matrix covers request drop, remote success plus response loss, timeout,
duplicate submission, crashes before and after submission, crash after remote
commit, old/new/both/neither acceptance, policy change, lockout risk, unavailable
old verification, and absent authoritative readback. The HTTP form adapter uses
new-acceptance evidence and never probes the old password. The LDAP-style
directory adapter uses expected-version evidence and idempotent operation
identity. They share one lifecycle but reach different safe terminal states.

All \RotationCases\ scenarios satisfy the commit-evidence and uncertainty-preservation
checks. \RotationCommitted\ reach committed because the configured adapter
predicate is satisfied; \RotationUnknown\ remain unknown,
\RotationPrepared\ remain prepared after a pre-submit crash, and
\RotationAborted\ abort after the directory model conclusively observes the base version.
The experiment uses high-fidelity local semantic models and real durable SQLite
reopen, but not a production HTTP application or OpenLDAP deployment.

\subsection{Bounded model checking}

The dedicated model represents four adapter evidence requirements, four remote
credential states, persistent preparation, crashes, unknown outcomes, and
sequential generations.

\begin{table}[t]
\centering
\caption{TLC results for the EPSCD rotation abstraction.}
\label{tab:tla}
\begin{tabular}{lr}
\toprule
Metric & Result \\
\midrule
Generated states & 2,877 \\
Distinct states & 1,069 \\
Search depth & 16 \\
Main invariant violations & 0 \\
Negative controls detected & 2/2 \\
\bottomrule
\end{tabular}

\end{table}

TLC exhaustively explores the configured finite state graph. Invariants cover
commit evidence, remote candidate acceptance, uncertainty reconstruction,
prepare-before-submit, no advance in unknown state, and sequential generation
uniqueness. The two negative controls enable a transition that commits on a
generic successful response and a transition that drops the candidate on
unknown outcome; the expected invariants are violated. The result validates
the abstract transition relation within its bounds and assumptions, not the
complete Rust implementation or an unbounded network.

\section{Security analysis}

\paragraph{Known credential pairs.}
Generation is public. Given $(g_i,\mathit{pwd}_{g_i})$, an observer can compute
$(g_i,\Rank_P(\mathit{pwd}_{g_i}))$, a known input/output pair for the
permutation. Security of unseen generations is the selected PRP's computational
claim, not information-theoretic secrecy or forward secrecy.

\paragraph{Credential-key and Root-Key compromise.}
Compromise of $K_{\mathrm{cred}}$ exposes its complete lineage. Replacing only
the policy epoch does not recover from that exposure because the epoch is public
context; a new credential salt and lineage identifier are required. Root-Key
compromise affects every derived credential context. Shamir resharing of the
same root changes share material but does not recover from root compromise.

\paragraph{Metadata theft.}
Policy, salt, generation, lineage, and adapter metadata are non-secret. Database
theft reveals account relationships and enables rollback attempts, but does not
alone yield the credential key under the root-secret assumption. Metadata
integrity and freshness remain necessary because changing generation or policy
context changes the reconstructed password.

\paragraph{Low-entropy accepted language.}
Uniformity cannot repair a small verifier domain. EPSCD reports exact effective
bits and can enforce a deployment threshold. The one-million-element FF1 floor
is an implementation/backend limit, not a sufficient password-strength rule.

\paragraph{Endpoint and adapter compromise.}
An endpoint compromised during legitimate reconstruction can observe plaintext
and key material. A compromised adapter can forge observations within its trust
boundary and defeat commit classification. Deployment must authenticate adapter
code and evidence sources; the state machine prevents accidental inference, not
malicious evidence by a trusted component.

\paragraph{Rollback and stale state.}
The journal prevents ordinary crash loss when the latest authenticated record is
available. A complete rollback of both local state and its freshness anchor is
outside that guarantee. An independent checkpoint can detect stale generation
and pending-operation metadata, but such a service is not required for the two
core constructions and is not evaluated as a contribution here.

\paragraph{Remote ambiguity and availability.}
Unknown outcome is a deliberate safety state. It may require operator action or
a later authoritative observation, and a target that exposes no safe evidence
may remain unresolved indefinitely. This sacrifices automatic liveness rather
than guessing a credential state.

\section{Limitations and deployment guidance}

The policy compiler supports a substantial bounded fragment, not arbitrary
natural-language policies or breach-list membership. Product-specific semantic
translation remains a trusted step. State construction can grow exponentially,
so fixed state, memory, and time budgets remain necessary even when all exact
translations complete within the final evaluation budget.

The FF1 dependency and cycle-walking wrapper are a prototype backend. The
1,024-walk cap creates explicit failure for rare inputs rather than total
derivation for every possible key/tweak/input. Timing varies with walk count and
the complete implementation is not constant time. A production deployment
would require independent cryptographic review and side-channel analysis.

The adapter evaluation does not include a real enterprise application,
OpenLDAP cluster, replication lag, CAPTCHA/MFA workflow, or measured human
reconciliation. The capability model is designed to express these differences,
but production evidence predicates must be validated per target. The formal
model is bounded and checks safety, not liveness.

Frequent forced password rotation is not itself recommended. EPSCD supports
necessary compromise response, inherited policy, and migration workflows; it
should not be used to justify needless rotation or to delay replacement of
password-only authentication.

\section{Conclusion}

EPSCD treats the complete bounded policy as the credential domain and treats
remote change success as an evidence decision. Exact counting and inverse
ranking expose the true domain; a context-separated exact-domain permutation
turns generation into a deterministic, same-lineage non-repeating sequence.
Durable preparation and adapter-specific evidence prevent local generation from
advancing on ambiguous transport outcomes while retaining both recovery paths.
The implementation, public-corpus run, fault matrix, and bounded model checking
support these scoped claims. The method composes established mathematical and
cryptographic tools into a legacy credential lifecycle; it does not claim new
primitives or remove the need to migrate legacy targets.

\section*{Artifact availability}

The artifact contains source, tests, the public-corpus translation provenance,
machine-readable results, generated tables, and a one-command reproduction
script. Identifying repository information can be replaced with an anonymous
artifact URL during double-blind review.

\bibliographystyle{elsarticle-num}
\bibliography{references}

\end{document}
